\section{温度与水分迁移模型}

\subsection{时变环境边界}
在附件1中截取 $0$ 至 $1800\,\mathrm{s}$ 的31个采样点，不作平滑或外推。对任一环境变量 $B\in\{\Tin,C_\infty\}$，在相邻采样时刻 $t_j,t_{j+1}$ 之间取
\begin{equation}
 B(t)=B_j+\frac{t-t_j}{t_{j+1}-t_j}(B_{j+1}-B_j),\qquad t_j\le t\le t_{j+1}.
 \label{eq:ambient}
\end{equation}
由此保留环境温湿度的实际变化，并避免预先设定恒温段。两端时刻的环境温度分别为 $28\,\celsius$ 和 $41.513\,\celsius$，环境水分数据分别为 $0.01963\,\kgkg$ 和 $0.03307\,\kgkg$。

\subsection{控制方程}
根据能量守恒与傅里叶定律，轴对称温度场满足
\begin{equation}
 \rho c_p\frac{\partial T}{\partial t}
 =\frac{1}{r}\frac{\partial}{\partial r}\left(rk\frac{\partial T}{\partial r}\right)
  +\frac{\partial}{\partial z}\left(k\frac{\partial T}{\partial z}\right).
 \label{eq:heat}
\end{equation}
在固定干物质体积密度的假设下，将水分守恒式中的该密度约去，得到
\begin{equation}
 \frac{\partial C}{\partial t}
 =\frac{1}{r}\frac{\partial}{\partial r}\left(rD(C)\frac{\partial C}{\partial r}\right)
  +\frac{\partial}{\partial z}\left(D(C)\frac{\partial C}{\partial z}\right),\quad
 D(C)=7\times10^{-9}\exp\left(-\frac{0.89}{C}\right).
 \label{eq:moisture}
\end{equation}
式\eqref{eq:moisture}必须保留散度形式。由于 $D$ 随 $C$ 变化，展开后有
\begin{equation}
 \nabla\!\cdot[D(C)\nabla C]
 =D(C)\Delta C+D'(C)|\nabla C|^2,\qquad
 D'(C)=\frac{0.89}{C^2}D(C).
 \label{eq:nonlinear}
\end{equation}
直接使用 $D(C)\Delta C$ 会遗漏非线性梯度项。第一问中热物性均为常数，且 $D$ 仅依赖 $C$，因此温度方程与水分方程可以分别求解。

\subsection{初始条件及边界条件}
初始状态为
\begin{equation}
 T(r,z,0)=T_0,\qquad C(r,z,0)=C_0.
 \label{eq:initial}
\end{equation}
在中心轴和中截面分别施加对称条件
\begin{equation}
 \left.\frac{\partial T}{\partial r}\right|_{r=0}
 =\left.\frac{\partial C}{\partial r}\right|_{r=0}=0,\qquad
 \left.\frac{\partial T}{\partial z}\right|_{z=0}
 =\left.\frac{\partial C}{\partial z}\right|_{z=0}=0.
 \label{eq:symmetry}
\end{equation}
令 $\Gamma_e=\{r=R\}\cup\{z=H\}$ 为半域的暴露边界，$\bm n$ 为外法向，$T_s,C_s$ 为表面未知量，则
\begin{equation}
 \begin{aligned}
 -k\nabla T\cdot\bm n&=h(T_s-\Tin),\\
 -D(C_s)\nabla C\cdot\bm n&=k_m(C_s-\Ceq),\qquad \Ceq=C_\infty.
 \end{aligned}
 \label{eq:robin}
\end{equation}
该符号约定下，表面较环境冷时热流向内，表面含水率较高时水分向外迁移。表面状态由式\eqref{eq:robin}求解，而非直接赋成环境值。侧面与端面交界处分别计算两个方向的通量。初始均匀含水率与瞬时失水边界在 $t=0$ 不相容，需保留题设初值并在时间离散中解析启动层。
